\documentclass[aspectratio=169,xcolor=dvipsnames]{beamer}
\usetheme{Berlin}

\usepackage[utf8]{inputenc}
\usepackage[T1]{fontenc}
\usepackage[spanish]{babel}
\usepackage{hyperref}
\usepackage{graphicx}
\usepackage{booktabs}
\usepackage{amsmath}
\usepackage{lettrine}
\setbeamertemplate{caption}[numbered]
\usepackage[dvipsnames,svgnames,x11names]{xcolor}
\usepackage{xurl}
\usepackage{algorithm}
\usepackage{algorithmicx}
\usepackage{algpseudocode}
\hypersetup{
    colorlinks=true,
    linkcolor=blue,
    filecolor=blue,
    urlcolor=blue,
    citecolor=blue,
}

%----------------------------------------------------------------------------------------
%   LISTINGS — PYTHON STYLE
%----------------------------------------------------------------------------------------
\usepackage{listings}
\definecolor{codegreen}{rgb}{0,0.6,0}
\definecolor{codegray}{rgb}{0.5,0.5,0.5}
\definecolor{codepurple}{rgb}{0.58,0,0.82}
\definecolor{backcolour}{rgb}{0.97,0.97,0.99}

\lstdefinestyle{PythonStyle}{
  language=Python,
  basicstyle=\ttfamily\footnotesize,
  keywordstyle=\color{blue}\bfseries,
  commentstyle=\color{codegreen},
  stringstyle=\color{violet},
  numberstyle=\tiny\color{codegray},
  breakatwhitespace=false,
  breaklines=true,
  captionpos=b,
  keepspaces=true,
  numbers=left,
  numbersep=5pt,
  showspaces=false,
  showstringspaces=false,
  showtabs=false,
  tabsize=2,
  frame=lines,
  framerule=0.4pt,
  backgroundcolor=\color{backcolour},
  extendedchars=true,
  inputencoding=utf8
}
\lstset{style=PythonStyle}

%----------------------------------------------------------------------------------------
%   TITLE PAGE
%----------------------------------------------------------------------------------------
\title{Estructuras de Selección Múltiple}
\subtitle{Materia: Algoritmos y Programación --- \texttt{match/case} en Python}

\author{Prof. D.Sc. BARSEKH-ONJI Aboud}
\institute{
    Facultad de Ingeniería \\
    Universidad Anáhuac México
}
\date{\today}

%----------------------------------------------------------------------------------------
%   AGENDA AUTOMÁTICA POR SECCIÓN
%----------------------------------------------------------------------------------------
\AtBeginSection[]{
  \begin{frame}{Agenda}
    \tableofcontents[currentsection]
  \end{frame}
}

\begin{document}

%----------------------------------------------------------------------------------------
\begin{frame}
    \titlepage
\end{frame}
%----------------------------------------------------------------------------------------

\begin{frame}{Agenda}
    \tableofcontents
\end{frame}

%================================================
\section{Motivación y Contexto}
%================================================

\begin{frame}{¿Por qué necesitamos selección múltiple?}
    \begin{columns}
        \column{0.5\textwidth}
        \textbf{Problema frecuente en programación:}
        \begin{itemize}
            \item Evaluar \textbf{una misma variable} contra muchos valores posibles.
            \item Con \texttt{if/elif} anidado el código crece y se vuelve difícil de leer.
            \item En lenguajes como C, Java o JavaScript existe la instrucción \texttt{switch/case}.
        \end{itemize}
        \column{0.5\textwidth}
        \begin{alertblock}{Python antes de la versión 3.10}
            Python \textbf{no} tenía \texttt{switch/case}. Los programadores usaban largas cadenas de \texttt{if/elif/else} o diccionarios de funciones.
        \end{alertblock}
        \vspace{0.3cm}
        \begin{block}{Python 3.10+ — Structural Pattern Matching}
            Se introduce la instrucción \texttt{match/case}, más poderosa que el \texttt{switch} clásico.
        \end{block}
    \end{columns}
\end{frame}

%================================================
\section{Sintaxis Básica de \texttt{match/case}}
%================================================

\begin{frame}[fragile]{Estructura General}
    \begin{columns}
        \column{0.45\textwidth}
        \textbf{Sintaxis:}
        \begin{lstlisting}[style=PythonStyle]
match variable:
    case valor1:
        # bloque 1
    case valor2:
        # bloque 2
    case _:
        # bloque por defecto
        \end{lstlisting}

        \column{0.55\textwidth}
        \begin{block}{Elementos clave}
            \begin{itemize}
                \item \texttt{match} — palabra clave, recibe la expresión a evaluar.
                \item \texttt{case valor} — compara con un patrón específico.
                \item \texttt{case \_} — comodín (\textit{wildcard}), equivale al \texttt{else}; se ejecuta si ningún caso anterior coincide.
                \item No es necesario \texttt{break} (a diferencia de C/Java).
                \item Disponible desde \textbf{Python 3.10}.
            \end{itemize}
        \end{block}
    \end{columns}
\end{frame}

\begin{frame}[fragile]{Ejemplo Mínimo — Día de la Semana}
\begin{lstlisting}[style=PythonStyle]
dia = int(input("Ingresa un numero de dia (1-7): "))

match dia:
    case 1:
        print("Lunes")
    case 2:
        print("Martes")
    case 3:
        print("Miercoles")
    case 4:
        print("Jueves")
    case 5:
        print("Viernes")
    case 6:
        print("Sabado")
    case 7:
        print("Domingo")
    case _:
        print("Numero invalido. Ingresa un valor entre 1 y 7.")
\end{lstlisting}

    \begin{exampleblock}{Nota}
        El caso \texttt{\_} actúa como valor por defecto — siempre se coloca al final.
    \end{exampleblock}
\end{frame}

%================================================
\section{Comparación con \texttt{if/elif}}
%================================================

\begin{frame}[fragile]{match/case vs.\ if/elif — Lado a Lado}
    \begin{columns}
        \column{0.5\textwidth}
        \textbf{Con \texttt{if/elif}:}
        \begin{lstlisting}[style=PythonStyle]
color = input("Color: ")

if color == "rojo":
    print("Stop")
elif color == "amarillo":
    print("Precaucion")
elif color == "verde":
    print("Avanza")
else:
    print("Color desconocido")
        \end{lstlisting}

        \column{0.5\textwidth}
        \textbf{Con \texttt{match/case}:}
        \begin{lstlisting}[style=PythonStyle]
color = input("Color: ")

match color:
    case "rojo":
        print("Stop")
    case "amarillo":
        print("Precaucion")
    case "verde":
        print("Avanza")
    case _:
        print("Color desconocido")
        \end{lstlisting}
    \end{columns}

    \vspace{0.2cm}
    \begin{alertblock}{Diferencia clave}
        \texttt{match/case} hace \textbf{comparación de patrones} (\textit{pattern matching}), no solo igualdad. Esto lo hace más expresivo que \texttt{if/elif}.
    \end{alertblock}
\end{frame}

\begin{frame}{Ventajas y Limitaciones}
    \begin{columns}
        \column{0.5\textwidth}
        \begin{block}{Ventajas de \texttt{match/case}}
            \begin{itemize}
                \item Código más \textbf{legible} y estructurado.
                \item Evita repetir la variable comparada.
                \item Soporta patrones complejos (tuplas, clases, guardas).
                \item Sin riesgo de olvidar \texttt{break}.
            \end{itemize}
        \end{block}

        \column{0.5\textwidth}
        \begin{alertblock}{Limitaciones}
            \begin{itemize}
                \item Solo disponible en \textbf{Python 3.10+}.
                \item No evalúa rangos directamente (se usa guarda \texttt{if}).
                \item No reemplaza al \texttt{if/elif} en comparaciones complejas con múltiples condiciones booleanas.
            \end{itemize}
        \end{alertblock}
    \end{columns}
\end{frame}

%================================================
\section{Patrones Avanzados}
%================================================

\begin{frame}[fragile]{Múltiples Valores en un Caso (\texttt{|} OR)}
\begin{lstlisting}[style=PythonStyle]
mes = int(input("Ingresa el numero de mes (1-12): "))

match mes:
    case 1 | 3 | 5 | 7 | 8 | 10 | 12:
        print(f"El mes {mes} tiene 31 dias.")
    case 4 | 6 | 9 | 11:
        print(f"El mes {mes} tiene 30 dias.")
    case 2:
        print("Febrero: 28 o 29 dias (bisiesto).")
    case _:
        print("Mes invalido.")
\end{lstlisting}

    \begin{block}{Operador \texttt{|} en \texttt{case}}
        El operador \texttt{|} permite agrupar varios valores en un mismo caso, equivalente a \texttt{elif x == 1 or x == 3 or ...}
    \end{block}
\end{frame}

\begin{frame}[fragile]{Guardas (\textit{Guards}) con \texttt{if}}
\begin{lstlisting}[style=PythonStyle]
nota = float(input("Ingresa la calificacion (0-10): "))

match nota:
    case n if n >= 9.0:
        print(f"Calificacion {n}: Excelente")
    case n if n >= 7.0:
        print(f"Calificacion {n}: Aprobado")
    case n if n >= 6.0:
        print(f"Calificacion {n}: Suficiente")
    case n if n >= 0.0:
        print(f"Calificacion {n}: Reprobado")
    case _:
        print("Calificacion fuera de rango.")
\end{lstlisting}

    \begin{exampleblock}{Guardas}
        Una \textit{guarda} (\texttt{case n if condicion}) permite agregar condiciones adicionales dentro del patrón. La variable \texttt{n} captura el valor de \texttt{nota}.
    \end{exampleblock}
\end{frame}

\begin{frame}[fragile]{Patrones con Tuplas}
\begin{lstlisting}[style=PythonStyle]
punto = (int(input("x: ")), int(input("y: ")))

match punto:
    case (0, 0):
        print("Origen")
    case (x, 0):
        print(f"Sobre el eje X en x={x}")
    case (0, y):
        print(f"Sobre el eje Y en y={y}")
    case (x, y):
        print(f"Punto en ({x}, {y})")
\end{lstlisting}

    \begin{block}{Desestructuración de tuplas}
        \texttt{match/case} puede \textbf{desestructurar} estructuras de datos. Las variables \texttt{x} e \texttt{y} dentro del \texttt{case} capturan automáticamente los valores del \textit{match}.
    \end{block}
\end{frame}

\begin{frame}[fragile]{Patrones con Clases y Objetos}
\begin{lstlisting}[style=PythonStyle]
from dataclasses import dataclass

@dataclass
class Figura:
    tipo: str
    lado: float = 0
    base: float = 0
    altura: float = 0

figura = Figura(tipo="triangulo", base=5.0, altura=3.0)

match figura:
    case Figura(tipo="cuadrado", lado=l):
        print(f"Area del cuadrado: {l**2:.2f}")
    case Figura(tipo="triangulo", base=b, altura=h):
        print(f"Area del triangulo: {b*h/2:.2f}")
    case Figura(tipo="rectangulo", base=b, altura=h):
        print(f"Area del rectangulo: {b*h:.2f}")
    case _:
        print("Figura no reconocida.")
\end{lstlisting}
\end{frame}

%================================================
\section{Ejemplo Integrador}
%================================================

\begin{frame}[fragile]{Calculadora con \texttt{match/case}}
\begin{lstlisting}[style=PythonStyle]
def calculadora(a, operacion, b):
    match operacion:
        case "+":
            return a + b
        case "-":
            return a - b
        case "*":
            return a * b
        case "/":
            if b != 0:
                return a / b
            else:
                return "Error: division entre cero"
        case "**":
            return a ** b
        case _:
            return f"Operacion '{operacion}' no reconocida."

print(calculadora(10, "+", 5))   # 15
print(calculadora(8,  "/", 0))   # Error
print(calculadora(2,  "**", 8))  # 256
\end{lstlisting}
\end{frame}

\begin{frame}[fragile]{Menú Interactivo de Aplicación}
\begin{lstlisting}[style=PythonStyle]
def menu():
    print("\n=== MENU PRINCIPAL ===")
    print("1. Ver lista de estudiantes")
    print("2. Agregar estudiante")
    print("3. Eliminar estudiante")
    print("4. Salir")

    opcion = input("Selecciona una opcion: ")

    match opcion:
        case "1":
            print("Mostrando lista...")
        case "2":
            nombre = input("Nombre del nuevo estudiante: ")
            print(f"Estudiante '{nombre}' agregado.")
        case "3":
            nombre = input("Nombre a eliminar: ")
            print(f"Estudiante '{nombre}' eliminado.")
        case "4":
            print("Saliendo del sistema.")
        case _:
            print("Opcion no valida. Intenta de nuevo.")

menu()
\end{lstlisting}
\end{frame}

%================================================
\section{Ejercicios para Casa}
%================================================

\begin{frame}{Ejercicio 1 — Conversor de Unidades}
    \begin{block}{Descripcion}
        Escribe un programa en Python que pida al usuario:
        \begin{enumerate}
            \item Un \textbf{valor numérico} (flotante).
            \item El \textbf{tipo de conversión} a realizar (elige una opción de menú).
        \end{enumerate}
    \end{block}

    \vspace{0.3cm}
    El programa debe usar \texttt{match/case} para manejar las siguientes conversiones:

    \begin{columns}
        \column{0.5\textwidth}
        \begin{itemize}
            \item \texttt{"km\_mi"} $\rightarrow$ km a millas \hfill ($\times\,0.621$)
            \item \texttt{"mi\_km"} $\rightarrow$ millas a km \hfill ($\times\,1.609$)
            \item \texttt{"kg\_lb"} $\rightarrow$ kg a libras \hfill ($\times\,2.205$)
        \end{itemize}
        \column{0.5\textwidth}
        \begin{itemize}
            \item \texttt{"lb\_kg"} $\rightarrow$ libras a kg \hfill ($\div\,2.205$)
            \item \texttt{"c\_f"}   $\rightarrow$ Celsius a Fahrenheit \hfill ($\times\frac{9}{5}+32$)
            \item \texttt{"f\_c"}   $\rightarrow$ Fahrenheit a Celsius \hfill ($(\text{val}-32)\times\frac{5}{9}$)
            \item Cualquier otro valor $\rightarrow$ mensaje de error.
        \end{itemize}
    \end{columns}

    \vspace{0.2cm}
    \begin{alertblock}{Requisito}
        Usar obligatoriamente \texttt{match/case}. Mostrar el resultado con 3 decimales usando f-string.
    \end{alertblock}
\end{frame}

\begin{frame}{Ejercicio 2 — Clasificador de Figuras Geométricas}
    \begin{block}{Descripcion}
        Escribe un programa que clasifique triángulos usando \textbf{patrones con guardas} en \texttt{match/case}.
    \end{block}

    \vspace{0.2cm}
    \textbf{El programa debe:}
    \begin{enumerate}
        \item Pedir al usuario los \textbf{tres lados} del triángulo ($a$, $b$, $c$) como flotantes.
        \item Verificar si los tres lados forman un triángulo válido: $a + b > c$, $a + c > b$, $b + c > a$.
        \item Usar \texttt{match/case} con \textbf{guardas} para clasificar:
    \end{enumerate}

    \vspace{0.2cm}
    \begin{columns}
        \column{0.5\textwidth}
        \begin{itemize}
            \item \textbf{Equilátero:} $a = b = c$
            \item \textbf{Isósceles:} exactamente dos lados iguales
            \item \textbf{Escaleno:} los tres lados diferentes
        \end{itemize}
        \column{0.5\textwidth}
        \begin{itemize}
            \item \textbf{Rectángulo:} cumple teorema de Pitágoras ($c^2 = a^2 + b^2$, con $c$ el mayor)
            \item \textbf{Triángulo inválido:} si no se cumple la desigualdad triangular
        \end{itemize}
    \end{columns}

    \vspace{0.2cm}
    \begin{alertblock}{Requisito extra}
        El programa debe también calcular e imprimir el \textbf{perímetro} y el \textbf{área} (fórmula de Herón) del triángulo si es válido.
    \end{alertblock}
\end{frame}

%================================================
\begin{frame}{Resumen}
    \begin{columns}
        \column{0.55\textwidth}
        \begin{block}{Lo que aprendimos hoy}
            \begin{itemize}
                \item \texttt{match/case} — disponible desde Python 3.10.
                \item Sintaxis básica y comodín \texttt{\_}.
                \item Operador \texttt{|} para múltiples valores.
                \item Guardas con \texttt{if} para rangos.
                \item Patrones con tuplas y clases.
                \item Comparación con \texttt{if/elif}.
            \end{itemize}
        \end{block}

        \column{0.45\textwidth}
        \begin{exampleblock}{Cuándo usar \texttt{match/case}}
            \begin{itemize}
                \item Menús de opciones.
                \item Procesamiento de comandos.
                \item Desestructuración de objetos.
                \item Código más legible que largas cadenas \texttt{elif}.
            \end{itemize}
        \end{exampleblock}

        \vspace{0.3cm}
        \begin{alertblock}{Entrega de ejercicios}
            Próxima clase. Archivos \texttt{.py} con comentarios explicativos.
        \end{alertblock}
    \end{columns}
\end{frame}

\end{document}
